\subsection{Implementation}
\label{rasc:sec:impl}

We implemented \sysname{} as a new component ($\sim$8.5K Python LOC) for Home Assistant (HA)~\cite{HAcore}, a widely used open-source smart-home platform. 
Deploying \sysname{} requires no changes to the HA core, integrations, or devices and their APIs. Users simply install \sysname{} and 
provide lightweight per-device-class mappings that bind action execution points (i.e., start, completion) to existing device state attributes. 
This lets \sysname{} infer action progress and schedule polls without modifying devices.
We also built a 500-LOC
HA simulator and adapted the \texttt{hass-virtual} third-party component~\cite{hassvirtual} for device simulation and Raspberry Pi-based emulation.

\begin{figure}[t]
\centering
\begin{minipage}{0.99\linewidth}
\begin{minted}[
  fontsize=\small,
  linenos,
  breaklines,
  frame=lines,
  numbersep=4pt,    % distance between numbers and code
  xleftmargin=9pt   % keeps numbers snug against the left frame
]{yaml}
- id: "1715802493830"
  alias: Heat when window starts closing and door is closed
  action:
    - parallel:
        - service: cover.close_cover
          target: { entity_id: cover.living_room_window }
        - service: cover.close_cover
          target: { entity_id: cover.balcony }
    - service: climate.set_temperature
      data: { temperature: 72 }
      depend_on: [start, complete]
      target: { entity_id: climate.main_thermostat }
\end{minted}
\caption{\bf \small Example RASC routine with \texttt{depend\_on} specification. 
{\it \texttt{desc}, \texttt{trigger}, \texttt{condition} omitted for brevity. The main thermostat is requested to heat to 72\textdegree F after the living room window has started closing \& the balcony door has closed. 
}}
\label{rasc:lst:script-example}
\end{minipage}
\end{figure}

\noindent \textbf{\abs{} API.}
We extend HA's action YAML with a single optional field, \texttt{depend\_on}. 
This field is an \emph{ordered list} of required progress events---one per parent, in the parents' order (e.g., \texttt{ack}, \texttt{start}, \texttt{complete})---enabling fine-grained sequencing.
An example routine with mixed dependencies appears in \Cref{rasc:lst:script-example}.
If \texttt{depend\_on} is omitted, we default to \texttt{complete}, preserving backward compatibility with HA's sequential semantics.
At runtime, the YAML is compiled into a dependency DAG; \scheduler{} monitors this DAG and triggers actions as soon as prerequisites are satisfied.

\noindent \textbf{Action length PDF.}
We obtain each action's probability distribution function by storing historical durations for (ack→start) and (start→complete). 
Actions are keyed by \(\langle\textit{device},\,\textit{action},\,\textit{transition}\rangle\). 
Consequently, \dynpoller{} requires a brief training phase (several trials) before it can return accurate results.

\noindent \textbf{Asynchronous poll placement computation.}
\Cref{rasc:algo:adaptive-polling-algorithm}'s compute time increases with the upper bound, which can be large for long actions. 
To avoid polling delays
in \sysname, we decouple computation from execution and make it asynchronous: upon an action's completion, poll placement is computed anew and stored for its next initiation. 

\noindent \textbf{Action status change notification.}
On each status change, \sysname{} publishes a \texttt{\abs\_\allowbreak RESPONSE} event on HA's event bus; applications subscribed to this topic receive the update.
